\documentclass[11pt]{article}
\usepackage[margin=1.1in]{geometry}
\usepackage{amsmath,amssymb,amsthm}
\usepackage{booktabs}
\usepackage[hidelinks]{hyperref}
\usepackage{microtype}

\newtheorem{definition}{Definition}[section]
\newtheorem{convention}[definition]{Convention}
\newcommand{\Q}{\mathbb{Q}}
\newcommand{\eng}{\texttt{certify.py}}

\title{Exponent Thresholds as Exact Linear-Programming Facets:\\
Certificates with Problem Instances as Data}
\author{Michael M.\ Ross\\
\small\href{mailto:michaelmross@cantab.net}{michaelmross@cantab.net}}
\date{\today\\[2pt]\small Draft}

\begin{document}
\maketitle

\begin{abstract}
Many published exponent thresholds in analytic number theory arise as the
optimum of a finite competition among estimates, each valid on a
rational-linear window of range exponents. We describe a small, general
tool that certifies such thresholds exactly. A problem instance is a data
file declaring block variables, slice parameters, decomposition pieces,
and tools encoded as negated admissibility conditions. The engine
enumerates failure polytopes and computes optima by exhaustive vertex
enumeration over $\Q$, with no floating point. A run is a certificate
only when its output is checked against independently declared expected
values taken from the source publication. We give two worked instances.
The first reproduces the threshold profile of a recent conditional
reduction for primes in the intervals $[4n^2-n,4n^2+n]$. The second
transcribes Rivat and Sargos's theorem on Piatetski-Shapiro primes and
certifies three published rationals from one atom table: the main
threshold $c<2817/2426$, the intermediate $c<441/380$, and the prior
record $c<15/13$, the latter two recovered as tool-removal ablations,
each with the witness collapsing onto the predicted facet. We state
precisely what such certificates do and do not establish, using the
companion paper of Rivat and Wu as the boundary case: its
exponential-sum layer is certifiable, while its sieve positivity layer,
which rests on numerical integration of Buchstab functions, is not.
\end{abstract}

\section{Introduction}

A recurring pattern in analytic number theory is the exponent record: a
rational number such as $2817/2426$ or $243/205$ marking how far a
method reaches, produced by balancing finitely many estimates over the
ranges of grouped variables. In the cleanest cases the record is
literally the optimum of a finite system of linear inequalities with
rational coefficients. The optimum is stated in the paper, and the
system is scattered through it: one window condition per lemma, one cap
per error term, a design choice fixing a parameter here, a combinatorial
identity imposing a ceiling there. Verifying the record means
reassembling that system and checking the arithmetic, which referees do
by hand, once.

This note describes a tool that makes the reassembled system a
first-class, executable object. The engine, \eng, is problem-agnostic
and fits in a few hundred lines. Everything problem-specific lives in a
JSON instance file: variables, parameters, decomposition pieces, and
tools, with every coefficient an exact rational. The engine enumerates
the polytopes on which every tool fails simultaneously and computes
exact optima by vertex enumeration over $\Q$. Its output is compared
against expected values declared in the instance and taken from the
source publication, and the run ends with a verdict. We call a passing
run an \emph{exponent certificate} (Definition~\ref{def:cert}). The
same machinery answers a second question at essentially no cost: given
a hypothetical new estimate, encoded as one more tool, does the
threshold move, and where does the binding constraint migrate? A
published record thus acquires an executable acceptance test for
candidate improvements.

Two instances accompany the tool. The first certifies the threshold
profile of a conditional reduction for primes in square-centered
intervals \cite{ross-reduction,ross-inventory}, reproducing the profile
$\zeta(\gamma)=(6-5\gamma)/4$, its witness, and a tool-ablation
byproduct. The second transcribes Rivat and Sargos
\cite{rivat-sargos} and is the better advertisement for the method: a
single atom table certifies three published rationals, two of them as
ablations, all attained on the predicted facet
(Section~\ref{sec:rs}). The companion paper of Rivat and Wu
\cite{rivat-wu} supplies the boundary of scope
(Section~\ref{sec:limits}).

\section{The certification framework}\label{sec:framework}

\begin{definition}[Instance]\label{def:instance}
An \emph{instance} consists of block variables
$x=(x_1,\dots,x_n)$, slice parameters $\theta=(\theta_1,\dots,\theta_k)$,
a designated objective coordinate $x_j$, a finite set of \emph{tools}, a
finite set of \emph{pieces}, a finite list of rational slice values, and
declared expected values. A \emph{constraint} is an inequality
$a\cdot x\le b_0+b\cdot\theta$ with $a\in\Q^n$, $b_0\in\Q$,
$b\in\Q^k$. A tool is the negation of an estimate's admissibility
condition, written in disjunctive normal form: a list of \emph{fail
options}, each a conjunction of constraints, the tool failing at $x$
when at least one option holds. A piece is a pair (region, applicable
tools), the region being a bounded rational-linear set of block
coordinates that the piece must cover.
\end{definition}

Fix a slice $\theta$. A \emph{total-failure witness} in a piece is a
point of its region at which every applicable tool fails. The engine
finds witnesses by distributing the disjunctions: for each choice of
one fail option per tool it forms the polytope cut out by the region
together with the chosen fail atoms, and it examines every vertex,
solving each $n$-subset of constraint rows exactly over $\Q$. Two
certified statement types are supported.

\emph{Minimization mode.} The certified object is
$\zeta(\theta)=\min x_j$ over all witnesses at slice $\theta$: the
first level at which total failure occurs. The expected value is a
closed form, affine in $\theta$ with an optional floor.

\emph{Coverage mode.} The certified object is witness existence per
slice. The expected value is a threshold $\theta^*$ together with the
assertion that a witness exists exactly for $\theta\le\theta^*$.

\begin{convention}[Facet convention]\label{conv:facet}
All inequalities are closed, and the $\varepsilon$, $\eta$, $\kappa$
margins of the source arguments are dropped. Certified thresholds are
therefore attained: at the critical slice the witness set collapses
onto the binding facet, and the engine reports the facet point. Reported
thresholds are the suprema of failing slices, matching the usual
convention that the analytic result holds on the open side.
\end{convention}

\begin{definition}[Exponent certificate]\label{def:cert}
An \emph{exponent certificate} for a published threshold is a run of
the engine on a declared instance whose computed values match, exactly
and in $\Q$, expected values declared in the instance and taken from
the source publication, for every slice and every declared ablation.
\end{definition}

The match requirement is what separates certification from
computation. A run with no declared expected values computes a profile
but certifies nothing, and the engine says so
(Section~\ref{sec:verdicts}).

\section{The soundness contract}\label{sec:verdicts}

The engine certifies the arithmetic of the encoded system exactly. It
does not certify that the encoding is faithful to the source paper.
Fidelity is an obligation on the instance author, discharged in the
instance file itself, which carries per-tool provenance to numbered
statements of the source and a documented list of encoding decisions.
The discipline follows the sandwich pattern of \cite{ross-inventory}.
Deliberate relaxations are recorded with their direction: omitting a
tool condition can only make failure easier, so certified thresholds
are conservative, and a hand-proved coverage argument closes the gap
on the relevant region. Design choices of the source, such as a
parameter the paper fixes rather than quantifies, are recorded as
substitutions, not variables. Window edges that the source covers by
construction, through $\eta$-shifted endpoints, carry no fail atom;
only genuine caps do.

A run ends in one of three verdicts. \textsc{Certificate passed} means
every computed value matches its declared expectation, which is the
reproduction of a published result. \textsc{Certificate failed} means
at least one mismatch, which is an audit finding: either a
transcription error, to be fixed and rerun, or a genuine discrepancy
with the source, to be investigated by hand. \textsc{Nothing
certified} means the instance declared no expected values, which is
the exploration regime: the acceptance-test workflow of
Section~\ref{sec:acceptance} produces predictions awaiting theorems,
not checks against them. Keeping the third verdict distinct prevents
an exploratory run from masquerading as a certificate.

The transcription of an instance from a paper is mechanical enough to
delegate, to a careful reader or to an AI assistant working under the
protocol distributed with the tool \cite{repo}, because the acceptance
test is self-checking. An unfaithful atom table has no mechanism for
reproducing the source's exact published rationals, let alone several
of them simultaneously under ablation. The judgment calls concentrate
in two places: distinguishing genuine caps from design choices, and
detecting where the source leaves the rational-linear regime. Those
calls are audited by a human against the paper. The arithmetic audits
itself.

\section{Instance A: a square-interval tail decomposition}
\label{sec:jn}

The first instance transcribes the constraint system of
\cite{ross-reduction}, whose companion inventory \cite{ross-inventory}
documents the encoding in full. The block variables are three exponents
$(\mu,\eta,\beta)$ of a bilinear decomposition, the slice parameter is
$\gamma$, and the toolkit comprises the Robert--Sargos trilinear
estimate \cite{robert-sargos} in four groupings, the Fouvry--Iwaniec
monomial estimate \cite{fouvry-iwaniec} in three, and an elementary
tool combining Kuzmin--Landau with the exponent pairs $(1/2,1/2)$ and
$(1/6,2/3)$. Two pieces carry the type I and type II ranges.

In minimization mode the engine reproduces the declared profile
$\zeta(\gamma)=(6-5\gamma)/4$ on $\gamma\in[1,6/5]$ at six slices,
exactly. At $\gamma=11/10$ the reported witness is
$(\mu,\eta,\beta)=(11/20,9/40,1/8)$, on the facet $4\beta+5\gamma=6$,
with the binding fail atom of each tool named in the output, matching
the inventory's hand-derived witness. A declared ablation removes the
three Fouvry--Iwaniec groupings and certifies that the profile is
unchanged, reproducing the inventory's byproduct that the coverage
argument does not need that estimate. The full run, profile plus
ablation, takes about four seconds.

\section{Instance B: Rivat--Sargos, three rationals from one table}
\label{sec:rs}

Rivat and Sargos \cite{rivat-sargos} prove
$\pi_c(x)\sim x/(c\log x)$ for $1<c<2817/2426$, where $\pi_c(x)$
counts $n\le x$ with $\lfloor n^c\rfloor$ prime. Writing
$\gamma=1/c$ and $\nu=\log_x N$ for the length of the distinguished
range, their argument runs Heath-Brown's identity \cite{heath-brown-id}
with $\ell=5$ and the design choice $\alpha=1-\gamma$, splitting the
required coverage into a type II window, a type I requirement, and a
residual type $\mathrm{I}'$ window $(2(1-\gamma),1/3]$ on which the
paper's new estimate acts. Every condition in the critical path is
rational-affine in $(\nu,\gamma)$. The type II tool of Heath-Brown
\cite{heath-brown-ps} caps at $\nu\le5\gamma-4$. The thesis type I
bound requires $\nu\ge3(1-\gamma)$. The new estimate, at
Th\'eor\`eme~1 strength, caps at
$\nu\le(939\gamma-761)/143$ together with $\nu\le\gamma-1/2$,
$\nu\le3\gamma-2$, and $\nu\le(5\gamma+9)/33$, and at Th\'eor\`eme~2
strength, before the improved fifth-derivative criterion of Sargos
\cite{sargos-crit}, at $\nu\le(147\gamma-119)/23$ and
$\nu\le(597\gamma-469)/133$ in place of the first and last caps.

The instance runs in coverage mode with both strengths present as
separate tools and two declared ablations. The certified results are
collected in Table~\ref{tab:rs}. In each row the threshold is an exact
rational, the witness at threshold sits at $\nu=1/3$ exactly, and the
threshold matches a rational stated in the source. The base toolkit
certifies Th\'eor\`eme~1. Removing the Th\'eor\`eme~1 tool certifies
Th\'eor\`eme~2. Removing both new tools leaves Heath-Brown covering
type $\mathrm{I}'$ up to $\nu=5\gamma-4$, and the certified threshold
is $\gamma^*=13/15$, that is $c<15/13$: precisely the prior record of
Liu and Rivat \cite{liu-rivat} that the paper identifies with the
trivial estimate of its central sum. Eighteen slice checks pass, and
the run takes under a second.

\begin{table}[ht]
\centering
\begin{tabular}{@{}llll@{}}
\toprule
Toolkit & Threshold $\gamma^*$ & Statement & Source claim\\
\midrule
Full & $2426/2817$ & $c<2817/2426$ & Th\'eor\`eme 1\\
Remove T1 tool & $380/441$ & $c<441/380$ & Th\'eor\`eme 2\\
Remove both new tools & $13/15$ & $c<15/13$ & Liu--Rivat baseline\\
\bottomrule
\end{tabular}
\caption{Certified thresholds for the Rivat--Sargos instance. Each is
an exact facet of the declared linear system, with the witness at
threshold at $\nu=1/3$. The two ablation rows are the paper's own
internal structure, recovered by tool removal.}
\label{tab:rs}
\end{table}

The table settles a small point of record. The secondary literature
transmits the Rivat--Sargos threshold in two versions, $2817/2426$ and
$2817/2425$. The paper states $2817/2426$, and the certificate
confirms that $2817/2426$ is the exact facet of the paper's own
system: at $\gamma=2426/2817$ the binding cap satisfies
$(939\gamma-761)/143=1/3$ identically.

\section{The acceptance-test workflow}\label{sec:acceptance}

Because tools are data, a candidate estimate is one more JSON stanza.
Adjoining it and rerunning answers, in seconds and in exact
arithmetic, whether the candidate moves the threshold and which
constraint binds next. In an experiment on Instance A, a fabricated
strengthening of one Robert--Sargos window lifted the profile at every
interior slice and migrated the binding piece from type II to the
elementary type I facet, locating the next obstruction without any
analysis. Such runs end in the third verdict, \textsc{nothing
certified}, and their outputs are predictions conditional on the
candidate becoming a theorem. The workflow inverts the usual economics
of checking: the months-long question ``would estimate $X$ improve
result $Y$'' becomes a transcription plus a run, and the published
record ships with a machine-checkable specification of what any
improvement must supply.

\section{What is not certified}\label{sec:limits}

Certification decides membership in the reachable set of a declared
toolkit, exactly, and decides nothing beyond it. No enumeration over
type I and type II atoms can manufacture a new atom, and for
prime-detection problems there is a principled obstruction behind
this: measures agreeing on all such data may disagree about primes,
which is the parity phenomenon, circumvented to date only by estimates
exploiting structure specific to a sequence
\cite{friedlander-iwaniec-parity}. A certificate that a threshold is
optimal for its toolkit is a certified impossibility for
rearrangement, and simultaneously a specification, through its active
facet, of the estimate that any crossing must supply.

The scope boundary is visible within a single problem and a single
year. Rivat and Wu \cite{rivat-wu} prove
$\pi_c(x)\gg x/(c\log x)$ for $c\le243/205$. Their exponential-sum
layer is certifiable in exactly the present sense: the windows of
their Corollaries 1 and 2, with endpoints $3-3\gamma$,
$(61\gamma-49)/11$, $(60-61\gamma)/11$, and $3\gamma-2$, are
rational-affine. Their sieve layer is not. The threshold $243/205$ is
the point at which a function $F(\lambda)$, a sum of two- to
four-dimensional integrals of the Buchstab function over simplex
domains, remains provably positive under symbolic and numerical
integration in Maple, with domains manually enlarged for
tractability. That is a different kind of computation, with different
verification needs, and an instance file must not silently span the
boundary between the two. The correct posture is to certify the linear
layer and state where it ends.

For the same reason, this tool is not formal verification. A
certificate checks that a declared linear system reproduces published
rationals under exact arithmetic. It does not verify the analysis
justifying the atoms, which is what a proof assistant would demand.
The contract here is deliberately lighter, and the per-atom provenance
required of instances is the bridge for anyone who wants to carry a
certified system further.

\section{Related work}\label{sec:related}

The nearest neighbors are the systematizations of exponent
information: the analytic number theory exponent database of Tao,
Trudgian, and Yang \cite{antedb} and the automated optimization of
exponent pairs \cite{trudgian-yang}, which derive records by searching
within known estimates, and the Polymath~8 project \cite{polymath8},
which optimized a fixed argument's numerology at scale. The present
tool is complementary to these: it certifies published thresholds as
exact facets of declared systems, recovers a paper's internal theorem
structure by ablation, and attaches an executable acceptance test to a
record. Instances are small data files, so certified systems could in
principle be contributed to, or checked against, a database of the
kind just cited.
%% TODO (MMR): tone and completeness of this section to taste; add or
%% trim citations, e.g. computer-assisted optimization in sieve theory,
%% verified numerics traditions, interval-arithmetic work.

\section{Availability}\label{sec:avail}

The engine, the schema with the transcription protocol, and both
instance files are available at \cite{repo}. The engine is a single
Python file with no dependencies beyond the standard library, using
\texttt{fractions.Fraction} throughout. Runs reported here complete in
seconds on commodity hardware. A frozen release with a DOI will
accompany this note.
%% TODO (MMR): create github.com/michaelmross/exponent-certificates,
%% mint Zenodo DOI, fill in below.

%% TODO (MMR): acknowledgments to taste, including AI-assistance
%% disclosure per your usual practice.

\begin{thebibliography}{99}

\bibitem{ross-reduction}
M.~M.~Ross,
\emph{A one-hypothesis reduction for primes in $[4n^2-n,\,4n^2+n]$},
Zenodo (2026). \url{https://doi.org/10.5281/zenodo.21220096}

\bibitem{ross-inventory}
M.~M.~Ross,
\emph{Constraint inventory for the $J_n$ tail decomposition},
repository document, \url{https://github.com/michaelmross/Legendre}.

\bibitem{rivat-sargos}
J.~Rivat and P.~Sargos,
\emph{Nombres premiers de la forme $\lfloor n^c\rfloor$},
Canad. J. Math. \textbf{53} (2001), 414--433.

\bibitem{rivat-wu}
J.~Rivat and J.~Wu,
\emph{Prime numbers of the form $[n^c]$},
Glasgow Math. J. \textbf{43} (2001), 237--254.

\bibitem{heath-brown-id}
D.~R.~Heath-Brown,
\emph{Prime numbers in short intervals and a generalized Vaughan
identity}, Canad. J. Math. \textbf{34} (1982), 1365--1377.

\bibitem{heath-brown-ps}
D.~R.~Heath-Brown,
\emph{The Pjatecki\u{\i}--\v{S}apiro prime number theorem},
J. Number Theory \textbf{16} (1983), 242--266.

\bibitem{sargos-crit}
P.~Sargos,
\emph{Un crit\`ere de la d\'eriv\'ee cinqui\`eme pour les sommes
d'exponentielles}, Bull. London Math. Soc. \textbf{32} (2000),
398--402.

\bibitem{liu-rivat}
H.-Q.~Liu and J.~Rivat,
\emph{On the Pjateckii-Shapiro prime number theorem},
Bull. London Math. Soc. \textbf{24} (1992), 143--147.

\bibitem{robert-sargos}
O.~Robert and P.~Sargos,
\emph{Three-dimensional exponential sums with monomials},
J. reine angew. Math. \textbf{591} (2006), 1--20.

\bibitem{fouvry-iwaniec}
E.~Fouvry and H.~Iwaniec,
\emph{Exponential sums with monomials},
J. Number Theory \textbf{33} (1989), 311--333.

\bibitem{friedlander-iwaniec-parity}
J.~Friedlander and H.~Iwaniec,
\emph{Using a parity-sensitive sieve to count prime values of a
polynomial}, Proc. Nat. Acad. Sci. USA \textbf{94} (1997), 1054--1058.

\bibitem{antedb}
T.~Tao, T.~Trudgian, and A.~Yang,
\emph{Database of known results on analytic number theory exponents},
\url{https://github.com/teorth/expdb}.
%% TODO (MMR): confirm preferred citation form for the ANTEDB.

\bibitem{trudgian-yang}
T.~S.~Trudgian and A.~Yang,
\emph{Toward optimal exponent pairs},
arXiv:2306.05599.

\bibitem{polymath8}
D.~H.~J.~Polymath,
\emph{New equidistribution estimates of Zhang type},
Algebra Number Theory \textbf{8} (2014), 2067--2199.

\bibitem{repo}
M.~M.~Ross,
\emph{exponent-certificates: exact certification of exponent
thresholds, instances as data},
\url{https://github.com/michaelmross/exponent-certificates}.

\end{thebibliography}

\end{document}
